\section{Scope, principal results, and conventions}
\label{scope:section}

This continuation starts at the explicit questions R1, R7, and R8 of
\emph{Harmonic Parity and Resolvent Identities}, preserved in the incoming
directory at the pinned revision \cite{IncomingParity,Repository}.
R1 asks for mixed even-tail moments, beginning with orders two and four.
R7 asks for a genuinely convergent generator with several free harmonic
orders and arbitrary zero blocks. R8 asks for a derivative transverse to
a harmonic index that is subsequently set to a nonpositive integer.
The corresponding results here are independent theorems with compatible
normalizations.

\subsection{What is established}

\begin{enumerate}[leftmargin=*,label=\textbf{\arabic*.}]
\item \textbf{All zero decorations with independent orders.}
Theorem~\ref{zbg:main} proves a single analytic formula for
\[
 \sum_{p_0,\ldots,p_k\geq0}
 \HH_{a,b}(0^{p_0},s_1,0^{p_1},\ldots,s_k,0^{p_k})
                     u_0^{p_0}\cdots u_k^{p_k}.
\]
Every $s_j$ is an independent complex variable. The series converges
locally normally for $|u_j|<1$, with a domain that does not shrink with
$k$. The proof gives all branch cancellations and does not infer an
analytic identity from agreement at integer endpoints.

\item \textbf{Transverse Gamma and Stieltjes gaps.}
Theorems~\ref{tr:gap} and~\ref{tr:integral} prove every order of
\[
 \left.\partial_u^r\HH_{a,b}(s,u,t)\right|_{u=-m},
 \qquad m\geq0,\quad r\geq1.
\]
The first new case is the log-Gamma gap
$\log\Gamma(z+1)-\log\Gamma(y)$ between two surviving lattice variables
$y>z$. There are explicit absolutely convergent Mellin integrals, finite
decorated-sum formulas, and a finite family sufficient after arbitrary
unmarked nonpositive indices have been removed. The simultaneous
Stieltjes completion at $u=1$ is included.

\item \textbf{Every mixed even-tail moment.}
Section~\ref{mt:section} evaluates all centered even moments of
$(\zeta(A)-H_n^{(A)})(\zeta(B)-H_n^{(B)})$ for even $A,B\geq2$.
The finite formula uses Bernoulli numbers and classical double zeta
values of odd weight; Euler's reduction removes the latter.
For each fixed pair $(A,B)$ a fixed finite span contains the entire
infinite sequence of moments. An analytic generating function and
finite ordinary-polylogarithm primitives complete the result.
This resolves the two-factor sector of R1; products with additional
harmonic factors are not covered by the same theorem.
\end{enumerate}

For a concrete illustration, put $x_n=n+\tfrac12$ and
$T_j(n)=\zeta(j,n+1)$. The unequal-order identities include
\begin{align}
 \sum_{n\geq0}T_2(n)T_4(n)
  &=\frac{15}{2}\zeta(5)-3\zeta(2)\zeta(3),\label{scope:first}\\
 \sum_{n\geq0}x_n^2T_2(n)T_4(n)
  &=\frac14\zeta(2)\zeta(3)+\frac16\zeta(2)
       +\frac12\zeta(3)-\frac58\zeta(5),\label{scope:second}\\
 \sum_{n\geq0}\left\{x_n^4T_2(n)T_4(n)-\frac13\right\}
  &=-\frac7{80}\zeta(2)\zeta(3)-\frac1{30}\zeta(2)
       -\frac14\zeta(3)+\frac7{32}\zeta(5)-\frac3{40}.
       \label{scope:third}
\end{align}
All three displayed series, including the brackets in the last one,
are ordinary convergent series. The third identity must retain its
subtraction. Since $T_2(n)=\psi'(n+1)$ and
$T_4(n)=\psi'''(n+1)/6$, multiplication by six gives corresponding
polygamma-product evaluations.

\subsection{Provenance and the precise claim of advancement}

The harmonic interpolant and its original geometric Mellin kernel are
established mathematics. Ihara, Nakamura, and Yamamoto prove their
entire interpolation and harmonic product law; their paper also
explains related earlier work of Komori
\cite[Proposition 3.3, Theorem 4.1, Corollary 4.2]{INY}.
The incoming independent-orders report identifies the same interpolant
and gives a relative endpoint normalization \cite{IncomingIndependent}.
The zero-index polynomial reductions and one-free-order generator are
already in \cite{IncomingParity}. We retain those definitions and give
the requested simultaneous analytic extension and transverse formulas.

Similarly, Bernoulli power sums, Hurwitz specializations, Gamma
asymptotics, and Euler's odd-weight double-zeta reduction are classical
\cite{DLMFHurwitz,DLMFGamma,DLMFBernoulli,Bradley2007}.
They are ingredients of the mixed-tail theorem, rather than separate
claims of new mathematics.

The present results are proposed additions relative to the inspected
repository. We make no assertion of literature-wide priority for every
specialization, no arithmetic-independence claim for a spanning set, and
no proof-assistant formalization claim. The original $S_6$ and revised
$S_8$ candidates are not proved or refuted here.

\subsection{Analytic and algebraic conventions}

All endpoint parameters $a,b$ lie in the right half-plane.
Principal logarithms define $(n+a)^{-s}$. Positive real integration
variables use their real logarithms. The Hurwitz zeta function is
normalized by
\begin{equation}\label{scope:stieltjes}
 \zeta(1+v,x)=\frac1v+
       \sum_{r=0}^{\infty}\frac{(-1)^r}{r!}\gamma_r(x)v^r,
 \qquad \gamma_0(x)=-\psi(x).
\end{equation}
A superscript $\zeta^{(r)}$ denotes differentiation in the first
argument. In particular \cite[\S25.11]{DLMFHurwitz},
\begin{equation}\label{scope:hurwitz}
 \zeta(-m,x)=-\frac{B_{m+1}(x)}{m+1},\qquad
 \zeta'(0,x)=\log\Gamma(x)-\frac12\log(2\pi).
\end{equation}
Bernoulli numbers mean $B_j=B_j(0)$, so $B_1=-1/2$ and $B_1(1)=1/2$.
The strict multiple-zeta convention is
\[
 \zeta(r,t)=\sum_{n>j\geq1}n^{-r}j^{-t}.
\]

For a word $w=(s_1,\ldots,s_k)$, define the meromorphic ordered tail
\[
 Z_c(w)=\sum_{n_1>\cdots>n_k\geq0}
                      \prod_j(n_j+c)^{-s_j},\qquad Z_c(\varnothing)=1.
\]
An initial domain with $\Re s_j>1$ suffices below. The relative
interpolant is the unique triangular solution of
\begin{equation}\label{scope:relative}
 Z_b(w)=\sum_{w=uv}Z_a(u)\HH_{a,b}(v),\qquad
 \HH_{a,b}(\varnothing)=1.
\end{equation}
The sum includes both empty prefix and empty suffix. Its depth-one value is
\begin{equation}\label{scope:h}
 h_{a,b}(s)=\zeta(s,b)-\zeta(s,a).
\end{equation}
The pole at $s=1$ cancels. Useful values are
\begin{align}
 h_{a,b}(1)&=\psi(a)-\psi(b),&
 h_{a,b}^{(r)}(1)&=(-1)^r\{\gamma_r(b)-\gamma_r(a)\},\label{scope:hjets}\\
 h_{a,b}'(0)&=\log\Gamma(b)-\log\Gamma(a),&
 h_{a,b}(0)&=a-b.\notag
\end{align}
At $a=b+N$, $N\in\N_0$, the relative word is the finite strict sum
over $N>n_1>\cdots>n_k\geq0$. For general endpoints it means
\eqref{scope:relative}, with the joint continuation proved again in
Section~\ref{zbg:section}. The finite endpoint property alone is not
our definition of continuation.

We use a cutoff constant only for a function with a specified asymptotic
expansion in powers and logarithms of that cutoff. It is the coefficient
of the monomial $N^0(\log N)^0$. Section~\ref{mt:section} proves when
that prescription agrees with the spectral value. This distinction
is essential: even an accidentally convergent paired boundary series
need not equal the analytic continuation from its initial half-plane.

